\centerline{\bf \Huge Слив региона Эйлера}

\begin{flushright}
        \scriptsize{Кто наберёт меньше 25 тот лох}
\end{flushright}

\problem{1}
Ненулевые вещественные числа $a$ и $b$ удовлетворяют условию

$$
\frac{a}{b+1}+\frac{b}{a+1}=1
$$


Какие значения может принимать выражение

$$
\frac{a}{b}+\frac{b}{a}-\frac{1}{a b} ?
$$

\problem{2}
Среди участников кругового шахматного турнира мальчиков втрое больше, чем девочек. Ничьих не было, а в сумме мальчики набрали столько же очков, сколько и девочки. Кто занял первое место: мальчик или девочка?

Напомним, что в круговом турнире каждый участник играет с любым другим участником ровно одну партию.

\problem{3}
Дан неравнобедренный треугольник $A B C$. На его сторонах $A B$ и $B C$ отмечены точки $K$ и $L$ такие, что прямые $K L$ и $A C$ параллельны. Отрезки $A L$ и $K C$ пересекаются в точке $S$. Известно, что $A K=A S$ и $K L=L C$. Докажите,что $A L=K B$.

\problem{4}
Клетки квадрата $20 \times 20$ раскрашены в несколько цветов. Известно, что каждая клетка граничит по стороне или углу хотя бы с двумя клетками того же цвета. В какое максимальное количество цветов может быть покрашена таблица?

\problem{5}
Докажите, что существует бесконечно много таких пар натуральных чисел ($m, n$), что $m$ и $n$ имеют одинаковые наборы простых делителей, и $m-1$ и $n-1$ также имеют одинаковые наборы простых делителей.


В остроугольном треугольнике $ABC$ точки $M$ и $N$ середины сторон $AB$ и $AC$ соответственно. $AD$ биссектриса угла $\angle BAC$. Касательные к описанной окружности треугольника $ABC$, проведённые в точках $B$ и $C$ пересекаются в точке $P$. Пусть прямая $PQ$ и биссектриса угла $\angle MON$ ($O$ центр описанной $(ABC)$) пересекаются в точке $Q$. Докажите, что $QA$ --- касательная к $(ABC)$.